%% ======================================================================
%% chap2/background_main.tex
%%
%% Chapter 2: Background and Related Work
%%
%% Surveys scholarly and industry publications across four themes:
%% blockchain-based content management, cross-chain interoperability,
%% decentralized monetization models, and the gaps that motivate the
%% research. Closes with Table 2.1 mapping eight identified gaps to
%% CCRMS contributions.
%%
%% Sections:
%%   2.1 Blockchain-Based Content Management (3 subsections)
%%   2.2 Cross-Chain Interoperability (5 subsections)
%%   2.3 Decentralized Monetization Models (6 subsections)
%%   2.4 Challenges and Gaps (Table 2.1 -- literature gaps and CCRMS
%%       contributions)
%%   2.5 Summary
%%
%% External labels referenced:
%%   subsec:ink-contract (5.2.4), subsec:royalty-propagation (5.2.3),
%%   subsec:comparative-analysis (7.1.7),
%%   subsec:decentralization-hhi (7.3.3),
%%   sec:ink-discontinuation (Appendix A.9)
%%
%% Citations: extensive -- approximately two-thirds of references.bib,
%% spanning DRM, cross-chain surveys, NFT/IP rights, monetization models,
%% smart contracts, the Polkadot ecosystem, and EU regulatory frameworks.
%% See in-line \citep{} and \citet{} commands for specific keys.
%% ======================================================================

\chapter{Background and Related Work}\label{chap:background}

This chapter surveys four themes: blockchain-based content management (Section~\ref{sec:blockchain-content-management}), cross-chain interoperability (Section~\ref{sec:cross-chain-interop}), decentralized monetization models (Section~\ref{sec:decentralized-monetization}), and the challenges and gaps that motivate the research (Section~\ref{sec:challenges-and-gaps}).

\section{Blockchain-Based Content Management}\label{sec:blockchain-content-management}

\subsection{Market Context}\label{subsec:market-context}

Three independent market research firms \citep{marketsandMarkets2025, mordorIntelligence2026, grandViewResearch2025} converge on estimates that the global DRM market is valued at approximately USD~6.7--6.9~billion in 2025--2026, with projections reaching USD~11--14~billion by 2030--2033, representing a Compound Annual Growth Rate (CAGR) of 10--11\%. This growth is primarily driven by streaming services, AI-enabled piracy, and regulatory pressures. The market is dominated by \textbf{content protection} solutions (such as Widevine, FairPlay, PlayReady), rather than \textbf{content monetization}.

\subsection{Foundational Work on Blockchain-Based DRM}\label{subsec:foundational-drm}

Blockchain-based Digital Rights Management (DRM) is an active research domain. \citet{garbaEtAl2021} incorporate digital watermarking into blockchain technology as a tamper-evident layer; \citet{patilEtAl2022} build on this by introducing the DRMChain hybrid encryption architecture; and \citet{yiEtAl2023} further extend the protection model with a redactable blockchain scheme that combines DRM with perceptual hashing, enabling selective content modification without invalidating ownership proofs. Both earlier studies prioritize protection mechanisms over the economics of content creators.

In the music industry, \citet{cirelloEtAl2023} derive a set of design principles for blockchain-based DRM systems through expert interviews and prototype testing, identifying transparency of rights ownership, automated royalty distribution, and programmable licensing as the three foundational requirements. \citet{cherdakovEtAl2023} and \citet{mendozaTelloEtAl2024} propose Web3 implementations that embed these principles in music-streaming smart contracts, demonstrating the feasibility of on-chain licensing but stopping short of cross-chain portability.

The most comprehensive recent work is \citet{madapatiPradhan2025}, whose Ethereum + IPFS framework achieves 409~TPS for registration; this comparison is misleading because their pipeline hashes off-chain and stores only results on-chain. \citet{wangEtAl2025} present a BLS threshold signature scheme achieving 40\,ms aggregation. Two Polkadot-specific studies \citep{xieLiu2022, xieTang2024a} propose conceptual cross-chain frameworks without working prototypes.

This literature establishes \textbf{blockchain-based DRM as a legitimate research domain}; its consistent limitation is that it emphasizes protection over monetization, a priority CCRMS reverses.

\subsection{NFTs and Intellectual Property Rights}\label{subsec:nfts-ip-rights}

Work on NFTs and intellectual property encompasses legal analysis \citep{uspto2024, davik2025, murray2022, eurojust2024} and standardization proposals \citep{reggianini2025, gattoEtAl2021, dreyerEtAl2021}. The most empirically significant study is \citet{darshanEtAl2025}: \textbf{80\% of mainstream NFTs fail to transfer copyright, only 12--30\% of listings include licensing information, and 50--60\% of secondary sales bypass creator royalties}. \citet{lahiriEtAl2026} confirm these findings via a meta-analysis of 125 NFT studies, identifying legal clarity around ownership transfers as the field's primary unsolved problem. \citet{cabotNadalEtAl2022} propose an ERC-721 extension enabling recipients to reject unwanted transfers, relevant to our permanent-ownership mode.

The most significant industry development is the February 2025 launch of \textbf{Story Protocol}, the first dedicated intellectual-property Layer-1 blockchain \citep{storyFoundation2025}. Built on the Cosmos SDK with EVM compatibility, it enables programmable IP registration, automated licensing, and on-chain royalty distribution. It is the most direct industry comparator to CCRMS: a single-chain L1 versus a cross-chain parachain.

\section{Cross-Chain Interoperability}\label{sec:cross-chain-interop}

\subsection{Polkadot Architecture}\label{subsec:polkadot-architecture}

Polkadot \citep{wood2016} is a heterogeneous multi-chain framework in which parachains share security through a central relay chain and communicate via XCM. The core primitives are parachains, the relay chain, validators, collators, and XCM. \citet{habermeierEtAl2021} and \citet{wood2022} provide the primary technical documentation for the XCM message format and multi-hop features that underpin our cross-chain implementation. The XCM protocol itself has received dedicated academic attention: \citet{morhacEtAl2023a} build ParaSpell, a developer SDK that abstracts XCM message construction across parachains, and \citet{morhacEtAl2023b} demonstrate XCMP interoperability at the protocol level. Together with \citet{morhacEtAl2025}, which proposes the Unispell universal adapter, they form the most coherent body of peer-reviewed academic work on XCM in practice.

Bridge-based systems such as Wormhole and LayerZero rely on external validator committees that could be compromised independently, enabling forged ownership proofs or hijacked royalty distributions. Polkadot's shared-security model eliminates this risk: all parachains use the same relay-chain validator set. Cross-chain rights operations use HRMP, which routes messages through the relay chain and inherits its security and finality.

\subsection{Cross-Chain Technology Surveys}\label{subsec:cross-chain-surveys}

The cross-chain literature matured rapidly from 2022 to 2025. \citet{renEtAl2023} classify interoperability solutions into five categories: sidechains, notary schemes, hashed time lock contracts, relays, and blockchain-agnostic protocols. \citet{maoEtAl2022} provide an earlier three-category taxonomy; \citet{bigiottiEtAl2025} extend both in ACM Computing Surveys. Two more recent surveys expand the field: \textbf{\citet{dengEtAl2025}} analyze more than 150 sources; \textbf{\citet{liEtAl2025a}} distinguish \textbf{asset interoperability} (token transfers) from \textbf{data interoperability} (cross-chain state verification). The latter distinction is directly relevant: CCRMS addresses both through XCM payments and \texttt{pallet-rights-verifier} Merkle proofs. \citet{kumarVEtAl2025} and \citet{caoEtAl2025} examine cross-chain efficiency and trustless interoperability, both aligned with CCRMS's design intent.
\subsection{Bridges and Snowbridge}\label{subsec:bridges-snowbridge}

Bridge architecture research includes \citet{yinEtAl2022} on multi-party computation, \citet{chenKEtAl2023} on reputation-based trust, and \citet{maricEtAl2025} on Isabelle/HOL formal verification of fail-safe bridges.

The most operationally significant development in blockchain infrastructure is \textbf{Snowbridge}, a trust-minimized bridge that enables interoperability between Polkadot and Ethereum. Snowbridge V2 was officially launched on the Polkadot mainnet on November 7, 2025, and includes support for arbitrary Ethereum contract calls from Polkadot, a 50\% fee reduction, and unordered message execution \citep{snowfork2025a}. The associated Polkassembly proposal \citep{snowfork2025b} documents more than a year of operational stability with no on-chain downtime. We position our Snowbridge integration as forward-compatible with V2, even though we developed the initial prototype against V1.

\subsection{Polkadot 2.0 and JAM}\label{subsec:polkadot-2-jam}

Two major ecosystem developments shaped our implementation timeline. \textbf{Polkadot 2.0} was comprehensively launched by October 2025, introducing Async Backing (eliminating relay-chain wait), Agile Coretime (on-demand block production), Elastic Scaling (multiple cores per parachain), and XCM v5 (multi-hop transactions) \citep{parityUpgrade2025, polkadotNetwork2025}. \textbf{JAM} (Join-Accumulate Machine, \citet{wood2024}) is expected to succeed the relay-chain model within the Polkadot ecosystem, with its testnet launched in January 2026 and mainnet scheduled for mid-to-late 2026. These developments confirm that the Polkadot architecture is \textbf{maturing rather than stagnating}.

\subsection{Cross-Chain Copyright Management as a Research Target}\label{subsec:cross-chain-copyright-research}

Numerous studies examine cross-chain copyright management as a research issue \citep{xieLiu2022, xieTang2024a, munsonEtAl2022, mareckova2024, oreroEtAl2023, buzu2021}; however, none have developed a functional end-to-end prototype. This gap, together with the architectural analysis above, motivates the selection of XCM and FRAME as the implementation substrate and justifies choosing Polkadot as the platform over single-chain or bridge-based alternatives.

XCM is preferred over Wormhole and LayerZero because it is native to the relay chain and introduces no additional trust assumptions. Its structured instruction format also encodes rights operations as typed, on-chain-interpretable messages rather than opaque byte payloads, preserving rights metadata across chain boundaries (gap~G3). FRAME is the implementation framework of choice over EVM and CosmWasm alternatives because its composable module system allows all three access modes to share a single pallet and access-check function, its weight benchmarking provides deterministic fee calculations, and its typed storage API exposes rights state as queryable on-chain objects via the RPC layer.

\section{Decentralized Monetization Models}\label{sec:decentralized-monetization}

\subsection{Creator Economy Context}\label{subsec:creator-economy-context}

\citet{communipass2026} forecasts the creator economy will exceed USD~280~billion by the end of 2026. The blockchain-specific segment is growing rapidly, from USD~499.52~million (2024) to USD~2.74~billion by 2030 \citep{iresearch3602025}, driven by creator dissatisfaction with centralized platforms and maturing blockchain infrastructure.

\subsection{DeFi and Web3 Foundations}\label{subsec:defi-web3}

The conceptual foundation for decentralized monetization lies in DeFi. \citet{gogelEtAl2021} present the canonical Wharton overview of the DeFi ecosystem, covering AMMs, lending, derivatives, and stablecoins; \citet{atzori2017} analyzes decentralized governance through a political economy lens; and \citet{busch2022} provides a technical overview of Web3 from the United States Congressional Research Service.

\subsection{Subscription, Pay-Per-View, and Purchase Models}\label{subsec:subscription-ppv-purchase}

Earlier work on blockchain subscription and pay-per-view models includes \citet{munson2019, khan2024, castweet2019, steem2025, banerjeeEtAl2020}. Most of these works examine blockchain incentives without addressing the micropayment problem. Two more recent works are particularly relevant: \textbf{\citet{osemwegie2025}} surveys decentralized media monetization (identifying the problems CCRMS targets but not proposing a working system), and \textbf{\citet{shah2025}} proposes a smart-contract framework for royalty distribution that is the closest to our royalty propagation design in the existing literature. CCRMS extends Shah's approach using a FRAME pallet with stronger atomicity guarantees. \citet{wamugo2024} provides empirical data on NFT adoption by media organizations in Kenya, demonstrating that blockchain monetization is under active consideration in emerging markets.

\subsection{Pay-Per-View and the Micropayment Problem}\label{subsec:ppv-micropayment}

Pay-per-view models have historically struggled on blockchain platforms because gas fees have often exceeded content costs. \citet{goyalEtAl2019} introduce Gringotts, a secure micropayment system, but do not address the core fee issue. \citet{chegenizadehEtAl2024} directly address the micropayment verification layer, proposing a multi-agent, privacy-preserving approach for resource-constrained offline devices that reduces per-transaction overhead through smart-contract batching; their work supports our strategic decision to model PPV as a counter rather than a sequence of individual on-chain transactions. \subsection{Smart Contracts as the Implementation Layer}\label{subsec:smart-contracts-impl}

Surveys of the smart contract literature \citep{ante2020, linEtAl2022, singhEtAl2024a, quanEtAl2024} establish the field as mature; Parity's ink! blog posts \citep{parityTechnologiesInk3_2022, parityTechnologiesInk_2022} were the primary technical references during implementation. The most consequential development is the \textbf{discontinuation of ink! in January 2026} \citep{inkAlliance2026}, discussed as gap G8 and in Appendix~\ref{sec:ink-discontinuation}.

\subsection{Royalty Automation}\label{subsec:royalty-automation}

Few peer-reviewed studies address on-chain royalty automation; \citet{shah2025} and \citet{reggianini2025} are the primary examples (discussed in Section~\ref{subsec:subscription-ppv-purchase}).

\section{Challenges and Gaps}\label{sec:challenges-and-gaps}

Table~\ref{tab:gaps} maps eight literature gaps to CCRMS contributions.

\begin{longtable}{@{}L{0.03\linewidth} L{0.27\linewidth} L{0.28\linewidth} L{0.36\linewidth}@{}}
\caption{Literature Gaps and CCRMS Contributions}\label{tab:gaps}\\
\toprule
\textbf{ID} & \textbf{Gap} & \textbf{Supporting literature} & \textbf{CCRMS contribution} \\
\midrule
\endfirsthead

\multicolumn{4}{l}{\emph{Table~\ref{tab:gaps} continued from previous page}}\\
\toprule
\textbf{ID} & \textbf{Gap} & \textbf{Supporting literature} & \textbf{CCRMS contribution} \\
\midrule
\endhead

\midrule
\multicolumn{4}{r}{\emph{Continued on next page}}\\
\endfoot

\bottomrule
\endlastfoot

G1 & Existing systems are protection-focused, not monetization-focused & \citet{garbaEtAl2021}; \citet{patilEtAl2022}; \citet{cirelloEtAl2023}; \citet{madapatiPradhan2025}; commercial DRM (Widevine, FairPlay, PlayReady) & Unified rights token with monetization as the primary design goal; protection becomes a secondary consequence of correct rights management. \\
\addlinespace
G2 & No native cross-chain recurring payments & Existing cross-chain protocols (XCM, Wormhole, LayerZero, Snowbridge) address one-time transfers; \citet{xieLiu2022, xieTang2024a} propose theoretical frameworks without recurring payments & On-chain \texttt{on\_initialize} block hook scheduler for automatic renewal. XCM v5 \texttt{Schedule} is not yet production-ready; we defer the implementation of true cross-chain scheduling to future work. \\
\addlinespace
G3 & Cross-chain protocols strip rights metadata on transfer & \citet{reggianini2025} addresses part of the standardization problem within European NFT licensing & \texttt{query\_rights\_metadata} extrinsic emitting a self-describing \texttt{RightsMetadata} event; metadata canonical on CCRMS, authenticated by \texttt{pallet-rights-verifier} Merkle proofs. \\
\addlinespace
G4 & Pay-per-view is economically unviable at micro-transaction scale on traditional blockchains & \citet{goyalEtAl2019}; \citet{chegenizadehEtAl2024}; \citet{khan2024}; x402 \citep{coinbaseX402_2025} & PPV reconceptualized as a counter: one payment per pack, cheap storage updates per view, with additive purchases (Appendix~\ref{sec:extrinsics-reference}). \\
\addlinespace
G5 & The decentralization--latency trade-off is poorly characterized quantitatively & Most works endorse decentralization qualitatively or measure only centralized performance; \citet{madapatiPradhan2025} is an exception but uses an off-chain pipeline; \citet{raoEtAl2024} provide a scalability review but no cross-paradigm comparison & Direct comparative benchmark against an Express.js + SQLite baseline (Section~\ref{subsec:comparative-analysis}) plus HHI calculation (Section~\ref{subsec:decentralization-hhi}) make the trade-off numerical. \\
\addlinespace
G6 & Regulatory uncertainty for tokenized rights (EU AI Act Aug 2025; MiCA full enforcement Jul 2026; EP report A10-0019/2026) & \citet{esma2025}; \citet{guadamuz2025} & Cryptographic provenance, immutable records, and automated royalty distribution align with emerging compliance infrastructure. MiCA classification of subscription tokens remains an open legal question. \\
\addlinespace
G7 & No existing system provides first-class native support for subscription, pay-per-view, and permanent ownership as three access modes within a single unified rights primitive. Story Protocol's PIL can approximate these commercial outcomes with separate license-token configurations; it does not treat renewal and expiry as native state transitions of a single content record. & Patreon (subscription only); OpenSea (ownership only); pay-per-view absent from most blockchain systems; Story Protocol (flexible licensing primitive, not unified rights); \citet{shah2025} addresses royalty automation as a single feature & Unified rights token: one \texttt{Content} record, three maps, one ordered \texttt{check\_access}, and renewal/expiry as first-class storage transitions, not as configurations of a generic licensing primitive. We do not claim that PIL cannot encode similar commercial outcomes. \\
\addlinespace
G8 & ink! discontinuation (January 2026) leaves Polkadot without a viable institutionally supported smart contract language; alternatives (\texttt{revive}, raw \texttt{pallet-revive}, \texttt{wrevive}) each carry trade-offs & Appendix~\ref{sec:ink-discontinuation} & Pallet-first architecture demonstrates that substantial functionality can be delivered without a smart contract language. The gap remains for systems whose business logic is better expressed in contracts. \\
\end{longtable}


We fully address gaps G1 and G7; we address G6 and G8 only partially because they depend on factors outside the prototype's control. The collective contribution is that a single working prototype demonstrates feasibility for a category of systems that the literature has addressed only partially.

\section{Summary}\label{sec:chapter-2-summary}

We surveyed the academic and industry literature on blockchain-based content management, cross-chain interoperability, and decentralized monetization. The foundational primitives are in place; the market for decentralized content monetization is real and expanding rapidly; and existing methodologies address the issue only partially. CCRMS bridges the eight identified gaps by deploying a unified rights record on a Polkadot parachain.

